\documentclass[11pt,a4paper]{article}
\usepackage[utf8]{inputenc}
\usepackage{amsmath,amssymb,amsthm}
\usepackage{booktabs}
\usepackage{float}
\usepackage{hyperref}
\usepackage{geometry}
\usepackage{orcidlink}
\geometry{margin=1in}
\setlength{\emergencystretch}{3em}

\newtheorem{theorem}{Theorem}
\newtheorem{proposition}{Proposition}
\newtheorem{lemma}{Lemma}
\newtheorem{conjecture}{Conjecture}
\newtheorem{remark}{Remark}
\newtheorem{definition}{Definition}

\title{A Quantitative Convergence Law for the\\
Connes--Consani--Moscovici Zeta Spectral Triple}

\author{Ronnie Andrews, Jr.\,\orcidlink{0009-0003-9724-3104}%
\thanks{Correspondence: \texttt{randrewsmath@gmail.com}}}

\date{June 2026}

\begin{document}
\maketitle

\begin{abstract}
The Connes--Consani--Moscovici (CCM) zeta spectral triple~\cite{CCM2025} and
its Connes--van Suijlekom refinement~\cite{CvS2025} produce finite real
symmetric matrices whose ground state yields an approximation $\nu_1$ to
the first
nontrivial Riemann zero $\gamma_1$ that sharpens as three parameters grow: a
prime cutoff $\lambda^2$, a basis size $N$, and the working precision $P$. The
rate of this convergence in the basis size is named the central open problem
in~\cite{CCM2025}. We characterize the convergence quantitatively, in two equivalent forms.
Intrinsically, the construction's relative error $|\nu_1-\gamma_1|/|\gamma_1|$
is a single two-channel quantity,
\[
  E(N,\lambda^2)=\varepsilon_{\mathrm{modes}}(N,\lambda^2)
                 +\varepsilon_{\mathrm{primes}}(\lambda^2),
\]
a mode channel vanishing as $N\to\infty$ and a prime channel vanishing as
$\lambda^2\to\infty$, so that $\nu_1\to\gamma_1$ is exactly $E\to0$ and forces
both limits. The two channels are the finite-$N$ transition and the extreme tail
of a single spectral edge (the lower edge of the equilibrium measure of the
archimedean field, whose bulk density is the Riemann--von Mangoldt law), so
$D_{\mathrm{nModes}}$ and $D_{\mathrm{Primes}}$ are one edge read at two scales. Observed at finite working precision $P$, the
same law is the matching-digit count
$D(\lambda^2,N,P)=-\log_{10}(|\nu_1-\gamma_1|/|\gamma_1|)$,
\[
  D(\lambda^2,N,P)=\min\!\bigl(D_{\mathrm{Wprec}}(P),\,
  D_{\mathrm{nModes}}(N,\lambda^2),\,D_{\mathrm{Primes}}(\lambda^2)\bigr),
\]
the minimum of three budgets --- precision, modes, and prime content --- of which
precision is the instrument resolution, absent from the intrinsic form. In digit
space the minimum is forced, not fitted ($D$ is a negative logarithm of a sum of
independent errors, so the largest binds), which is why earlier single-expression
digit fits fail. We prove that
convergence to $\gamma_1$ requires all three budgets to diverge: at a fixed
cutoff the accuracy is capped at the prime ceiling
$D_{\mathrm{Primes}}(\lambda^2)$, so no basis size or precision reaches the zero
and the cutoff must grow. This reduces convergence to the growth of the prime
ceiling, whose leading coefficient is derived, parameter-free, from the prolate
spheroidal eigenvalue asymptotic: $D_{\mathrm{Primes}}\sim(4\pi/\ln 10)\lambda^2$. The
mode budget
--- the basis-size convergence rate --- is the logarithm of a stretched-exponential
convergence $\exp(-aN^{q})$ whose exponent $q(\lambda^2)$ is the analyticity
class of the eigenvector, derived from the archimedean Gamma factor with a
prime correction from Mertens' theorem; a by-product is that the
eigenvalue-space Galerkin exponent measured by Groskin~\cite{Groskin2026} is a
finite-window slice of the same curve. We state each component with an explicit
rigor level, separating
what is proved versus what is derived and empirically validated up to
$\lambda^2=200$. The results have been communicated to
co-authors of~\cite{CCM2025}.
\end{abstract}

%==============================================================================
\section{Introduction}\label{sec:intro}
%==============================================================================

The Connes--Consani--Moscovici (CCM) construction~\cite{CCM2025}, building on
the Connes--van Suijlekom truncated Weil form~\cite{CvS2025}, produces a family
of finite self-adjoint operators $D_{\log}^{(\lambda,N)}$ whose spectra
approximate the imaginary parts $\gamma_k$ of the nontrivial Riemann zeta
zeros to remarkable numerical accuracy --- $55$ digits at the headline
configuration of~\cite{CCM2025}, which we extended to $999$ digits in an
independent reproduction~\cite{AndrewsA}. The construction has two integer-like
control parameters --- a prime cutoff $\lambda^2$ (the primes $p\le\lambda^2$
enter the Weil form) and a basis size $N$ (the matrix is $(2N{+}1)\times(2N{+}1)$)
--- together with the working precision $P$ in decimal digits. The rate at
which the eigenvalue converges in the basis size $N$ is named the central open
problem in~\cite{CCM2025}; Groskin~\cite{Groskin2026}
measured a Galerkin convergence exponent empirically but disclaims deriving it,
and \'Sliwi\'nski~\cite{Sliwinski2026} proves an inverse-logarithmic lower bound
on the mean spectral error. We previously
tested two proposed rate laws against high-precision data and found them
unsupported~\cite{AndrewsB}.

This paper characterizes that convergence explicitly.
\begin{definition}
$D(\lambda^2,N,P) := -\log_{10}\bigl(|\nu_1-\gamma_1|/|\gamma_1|\bigr)$ is the
\emph{matching-digit count}: the number of correct leading decimal digits of the
first Riemann zero $\gamma_1$ reproduced by $\nu_1$, the first ordinate produced
by the CCM construction, computed at working precision $P$ via MPFR/GMP. The reference $\gamma_1$ is
computed to $2500$ significant digits by rigorous interval arithmetic
(Arb~\cite{Arb2017}), far exceeding the largest match reported here
($\approx1070$ digits at $\lambda^2=200$), so the reported counts are genuine and
not reference-limited; its leading digits agree with the standard tabulations
of~\cite{Odlyzko}.
\end{definition}

\noindent Our central result, the \emph{Andrews CCM Convergence Law}, has two
equivalent forms. Its intrinsic (spectral) form is the two-channel error
\begin{equation}\label{eq:main-intrinsic}
  \boxed{\;E(N,\lambda^2)=\varepsilon_{\mathrm{modes}}(N,\lambda^2)
         +\varepsilon_{\mathrm{primes}}(\lambda^2)\;}
\end{equation}
whose two channels are the finite-$N$ transition and the extreme tail of a single
prolate-type spectral edge (Section~\ref{ssec:twoforms}); convergence
$\nu_1\to\gamma_1$ is exactly $E\to0$. Its finite-precision measurement is the
matching-digit count
\begin{equation}\label{eq:main}
  \boxed{\;D(\lambda^2,N,P)=\min\!\bigl(D_{\mathrm{Wprec}}(P),\,
  D_{\mathrm{nModes}}(N,\lambda^2),\,D_{\mathrm{Primes}}(\lambda^2)\bigr)\;}
\end{equation}
in which $D_{\mathrm{Wprec}}$ is the instrument resolution alone. Since $D$ is the
relative error on a logarithmic scale, the law is a quantitative convergence
statement: $D\to\infty$ is precisely $\nu_1\to\gamma_1$, and the three budgets
are the independent mechanisms that drive (or stall) that convergence
(Prop.~\ref{prop:limit}). Five points
distinguish it from prior empirical fits:
\begin{enumerate}
\item \textbf{The minimum is forced, not fitted} (Section~\ref{sec:skeleton}).
  In digit space $D$ is a negative logarithm of a sum of independent positive
  errors, so the largest dominates and $D$ equals the minimum of the three
  per-source budgets up to an $O(1)$ term --- the structural reason earlier
  single-formula digit fits failed. The precision budget is measurement, not
  mathematics: removing it leaves the single intrinsic
  edge~\eqref{eq:main-intrinsic}, whose two channels (modes, primes) are the only
  mathematical content.
\item \textbf{Each budget is analyzed from the construction}, not guessed. The
  precision budget follows from elementary round-off analysis; the prime
  ceiling from the prolate eigenvalue asymptotic; and the mode budget from the
  analyticity of the Weil-form eigenvector.
\item \textbf{Every claim carries an explicit rigor level} (Table~\ref{tab:tiers}),
  separating the rigorously proved parts from those derived and validated
  against high-precision computation. We do not label the whole law a
  theorem; we state precisely which pieces are.
\item \textbf{It yields a convergence criterion} (Prop.~\ref{prop:limit}): the
  prime cutoff is an unbeatable ceiling --- at a fixed $\lambda^2$ no basis size
  or precision reaches the zero, so convergence to $\gamma_1$ forces
  $\lambda^2\to\infty$ --- which turns the qualitative convergence
  of~\cite{CCM2025} into an explicit, conditional rate.
\item \textbf{No prior unified law, to our knowledge.} Earlier attempts are
  single closed-form fits of the digit count, and they fail: we tested two
  proposed rate laws against high-precision data and found them
  unsupported~\cite{AndrewsB}, and Groskin~\cite{Groskin2026} likewise reports
  that no single smooth-convergence law fits his data and declines to assert a
  rate, while \'Sliwi\'nski~\cite{Sliwinski2026} establishes only a lower bound
  on the mean error. The unified characterization given here --- the intrinsic
  two-channel edge~\eqref{eq:main-intrinsic} and its finite-precision
  measurement~\eqref{eq:main} --- is, to the best of the author's knowledge, the
  first working closed-form convergence law for the construction. Its unity is
  the point: one spectral edge (Section~\ref{sec:equilibrium}) underlies both the
  mode ramp and the prime ceiling, which is why a single expression can hold
  where earlier per-fits could not.
\end{enumerate}

The deepest new content is the mode budget $D_{\mathrm{nModes}}$
(Section~\ref{sec:nmodes}): we show the matching-digit ramp is the logarithm of
a stretched-exponential convergence, identify its exponent with the analyticity
class of the eigenvector, and derive that exponent end to end. A by-product is
that the eigenvalue-space Galerkin exponent measured by Groskin~\cite{Groskin2026}
is a finite-window slice of the same curve. That the mode ramp and the prime
ceiling are not two mechanisms but one spectral edge at two scales --- the lower
edge of an equilibrium measure whose bulk is the Riemann--von Mangoldt density
--- is the subject of Section~\ref{sec:equilibrium}.

%==============================================================================
\section{The construction and the object $D$}\label{sec:background}
%==============================================================================

We recall the construction of~\cite{CCM2025,CvS2025}. Let $\lambda>1$ and
$L=2\ln\lambda=\ln\lambda^2$. On $\mathcal H_\lambda=L^2([\lambda^{-1},\lambda],
du/u)$ the functions $V_n(u)=L^{-1/2}e^{2\pi i n\ln(\lambda u)/L}$ form an
orthonormal basis. The Weil quadratic form $QW_\lambda$, restricted to
$\mathrm{span}\{V_n:|n|\le N\}$, is a real symmetric $(2N{+}1)\times(2N{+}1)$
matrix whose archimedean part carries the local factor at infinity and whose
arithmetic part is a finite sum over prime powers $p^j\le\lambda^2$. Its
smallest eigenvalue $\varepsilon_N$ has an even eigenvector $\xi$, and the
positive zeros of the associated Fourier--Mellin function $\widehat\xi$
approximate the ordinates $\gamma_k$~\cite{CvS2025}.

Two facts about $\widehat\xi$ are used below and are due to~\cite{CCM2025,CvS2025}:
$\widehat\xi$ is entire, and (in the critical case) all its zeros are real and
coincide with the spectrum. The eigenvector $\xi$ is even and, in the regime we
study, the smallest eigenvalue is simple --- the even-symmetry verified
numerically across $38$ configurations spanning $\lambda^2=13$--$1000$
in~\cite{AndrewsA}.

The object of study, $D(\lambda^2,N,P)$, is a digit count: $D=-\log_{10}
(\text{relative error})$. The relative error itself,
$E=|\nu_1-\gamma_1|/|\gamma_1|$, is the intrinsic object of the
law~\eqref{eq:main-intrinsic}; $D$ is its measurement in digits, and the
arithmetic of that measurement is what makes the minimum structure
of~\eqref{eq:main} inevitable.

%==============================================================================
\section{The three-budget skeleton}\label{sec:skeleton}
%==============================================================================

The computed approximation $\nu_1$ carries error from three independent sources:
finite working precision (round-off), finite basis size (mode truncation), and
finite prime cutoff (the intrinsic ceiling at $N,P\to\infty$):
\[
  \text{error}\approx \varepsilon_{\mathrm{round}}(P)
  +\varepsilon_{\mathrm{modes}}(N,\lambda^2)+\varepsilon_{\mathrm{primes}}(\lambda^2).
\]
They are independent in that any one can be made dominant by taking the other
two parameters large.

\begin{proposition}[The minimum is forced]\label{prop:min}
For positive errors $\varepsilon_1,\varepsilon_2,\varepsilon_3$ with
$\Sigma=\sum_i\varepsilon_i$,
\[
  -\log_{10}\Sigma=-\log_{10}\bigl(\max_i\varepsilon_i\bigr)+O(1)
  =\min_i\bigl(-\log_{10}\varepsilon_i\bigr)+O(1),
\]
since $\max_i\varepsilon_i\le\Sigma\le 3\max_i\varepsilon_i$ and
$\log_{10}3<0.48$. Hence $D=\min(D_{\mathrm{Wprec}},D_{\mathrm{nModes}},
D_{\mathrm{Primes}})+O(1)$.
\end{proposition}

This is a rigorous consequence of two facts only: that $D$ is a digit count,
and that the three sources are independent. No fitting is involved. By the
bound, the minimum agrees with $D$ to within $\log_{10}3<0.48$ of a digit, so the
law's content is that the \emph{binding mechanism} changes across the parameter
space --- round-off, mode truncation, then prime content --- not that $D$ is
literally non-smooth (it is smooth: a negative logarithm of a sum).

The three sources are independent because they answer to different parameters
--- round-off to $P$, truncation to $N$, and prime content to $\lambda^2$ --- so
any one can be driven below the others at will. The $O(1)$ term is at most
$\log_{10}3<0.48$ of a digit, and arises only in the narrow crossover where two
budgets are comparable; away from those seams a single budget binds and $D$
equals it to within rounding. This is why, in digit space, $D$ is a minimum rather than a sum or product, and
why global single-expression fits of the digit count fail in practice: a single
elementary term cannot simultaneously hold a hard precision wall, a prime
ceiling, and a $\lambda^2$-coupled mode ramp. In the intrinsic form the error is,
by contrast, literally a sum, $E=\varepsilon_{\mathrm{modes}}
+\varepsilon_{\mathrm{primes}}$~\eqref{eq:main-intrinsic}; the digit-space minimum
is what the negative logarithm of that sum becomes, the largest term binding. The junctions differ in character.
The precision junction is a hard wall --- round-off is a clamp --- so its corner
is sharp. The mode$\to$prime junction, by contrast, is a wide, smooth saturation
(the $\tanh$ of Section~\ref{sec:nmodes}): $D_{\mathrm{nModes}}$ is the
mode-limited digit count, which rises along its pre-saturation ramp and is capped
by the ceiling, $D_{\mathrm{nModes}}\le D_{\mathrm{Primes}}$, with equality as
$N\to\infty$.

Each budget is isolated experimentally by \emph{starvation}: to read $D_i$
directly, set the other two parameters large enough that their budgets exceed
$D_i$. We use this throughout Sections~\ref{sec:wprec}--\ref{sec:primes}.

\subsection{The two forms of the law}\label{ssec:twoforms}

Of the three budgets, only two --- mode truncation and prime cutoff --- are
properties of the finite operator; round-off is a property of the computation.
Separating them gives the two forms stated in Section~\ref{sec:intro}.

\smallskip
\noindent\emph{Intrinsic (spectral) form.} The mathematical content is the
two-channel error~\eqref{eq:main-intrinsic}, $E=\varepsilon_{\mathrm{modes}}
+\varepsilon_{\mathrm{primes}}$, with $\varepsilon_{\mathrm{modes}}\sim
\exp(-aN^{q})\to0$ as $N\to\infty$ (Section~\ref{sec:nmodes}) and
$\varepsilon_{\mathrm{primes}}\to0$ as $\lambda^2\to\infty$
(Section~\ref{sec:primes}). It makes no reference to precision, and
$\nu_1\to\gamma_1$ is exactly $E\to0$, which requires both channels to vanish.

\smallskip
\noindent\emph{Observed (digit) form.} A computation at working precision $P$
resolves $E$ only down to its round-off floor, so the measured matching-digit
count is
\begin{equation}\label{eq:observed}
  D(\lambda^2,N,P)=\min\!\bigl(D_{\mathrm{Wprec}}(P),\,-\log_{10}E(N,\lambda^2)\bigr)
   =\min\!\bigl(D_{\mathrm{Wprec}},D_{\mathrm{nModes}},D_{\mathrm{Primes}}\bigr)+O(1),
\end{equation}
the second equality by Proposition~\ref{prop:min}. Here $D_{\mathrm{Wprec}}(P)$
is the instrument resolution: it is no property of $\nu_1$, $\gamma_1$, or the
operator, and as $P\to\infty$ it disappears, returning the intrinsic form in
digit units. Thus the boxed law~\eqref{eq:main} is the finite-precision
\emph{measurement} of the intrinsic law~\eqref{eq:main-intrinsic}.

\smallskip
\noindent\emph{One edge, not two budgets.} The mode and prime channels are the
two ends of a single spectral edge. Dividing out the ceiling
(Section~\ref{sec:nmodes}) writes $D_{\mathrm{nModes}}=D_{\mathrm{Primes}}\cdot g$
with $g=\tanh\!\bigl((c\,N/(\lambda^2\ln\lambda^2))^{q}\bigr)\uparrow1$ as
$N\to\infty$, so $D_{\mathrm{Primes}}=\lim_{N\to\infty}D_{\mathrm{nModes}}$ is the
$N\to\infty$ ceiling of the mode ramp, not a separate term; equivalently the
digit law is a single edge function $D_{\mathrm{edge}}(N,\lambda^2)=
D_{\mathrm{Primes}}\cdot g$ capped by precision,
$D=\min(D_{\mathrm{Wprec}},D_{\mathrm{edge}})+O(1)$. Correspondingly the intrinsic
error~\eqref{eq:main-intrinsic} is one edge: $\varepsilon_{\mathrm{modes}}$ is its
finite-$N$ transition (the Landau--Widom plunge, Section~\ref{sec:nmodes}) and
$\varepsilon_{\mathrm{primes}}$ is its extreme tail (the prolate asymptotic,
Section~\ref{sec:primes}). That these are literally the transition and the tail
of one band edge --- the lower edge of the archimedean equilibrium measure ---
is established in Section~\ref{sec:equilibrium}.

%==============================================================================
\section{The precision budget $D_{\mathrm{Wprec}}$}\label{sec:wprec}
%==============================================================================

Finite working precision $P$ (decimal digits) introduces relative round-off
$\sim 10^{-P}$, amplified by the eigenproblem's condition number $\kappa$ (set
by the gap of $\varepsilon_N$ to the rest of the spectrum and independent of
$P$). Hence
\begin{equation}\label{eq:wprec}
  D_{\mathrm{Wprec}}(P)=P-\log_{10}\kappa = P+O(1),
\end{equation}
with leading coefficient exactly $1$. The $O(1)$ term is implementation
dependent: the toolkit used here allocates $\lceil P\cdot\ln10/\ln2\rceil+16$
working bits, the $16$ guard bits contributing a small positive offset. Isolating
this budget (large $\lambda^2$ and $N$) and sweeping $P\in\{100,200,300,400,700,
1000\}$ gives a measured slope $\partial D/\partial P=1.000\pm0.001$. The exact
offset is not a property of the construction and we do not assign it physical
meaning. \emph{Status: T1; the coefficient $1$ follows from elementary numerical
analysis.}

\subsection{Symmetry, positivity, and the precision floor}\label{ssec:evenness}

A question the precision floor settles --- and one easy to answer wrongly
without it --- is whether the even symmetry of the ground eigenvector, and the
positivity of $\varepsilon_N$, are genuine features of the operator or merely of
the regime above the floor. The CCM construction~\cite{CCM2025} posits a simple
smallest eigenvalue with an \emph{even} eigenvector (its Step~1) and enforces
this with a forced-even projection. The reproduction of~\cite{AndrewsA} tested
that projection directly: across $38$ configurations ($\lambda^2=13$--$1000$,
$N=10$--$800$, spanning both sides of the precision floor) with the projection
\emph{disabled}, the natural smallest eigenvector converged to the very vector
the projection would have produced --- to \emph{bit-identical} eigenvalues ---
at every configuration above the floor. The forced-even projection is thus
empirically a no-op above the floor: even symmetry is a property the operator
exhibits unbidden, not an artifact of the projection or of the above-floor
regime.

The qualifier ``above the floor'' is essential, and it is precisely where the
budget of this section bites. When the target eigenvalue lies below
$D_{\mathrm{Wprec}}$ --- smaller than the round-off that the working precision
can resolve --- the computed eigenvector is dominated by numerical noise: it can
present spurious mixed symmetry, and the smallest computed eigenvalue can carry
a spurious sign. These are artifacts of precision starvation, not properties of
the operator; raising $P$ past the floor dissolves them, and the natural
eigenvector is again even, into the deep near-zero spectrum.

This even, positive ground state was documented in our reproduction~\cite{AndrewsA}.
In subsequent correspondence, and in a revision of his own paper~\cite{Groskin2026},
Groskin independently corroborated it by a different route: the sign changes his
separate quadrature implementation had shown among the deepest eigenvalues at a
finite archimedean cutoff $T$ were identified as a finite-$T$ artifact that
vanishes as $T$ is increased, consistent with the naturally even, positive ground
state of~\cite{AndrewsA}. Working precision $P$ and archimedean cutoff $T$ are
therefore independent controls: a deep-spectrum value is trustworthy only when it
is stable across two settings of \emph{each}, not when it merely holds steady in
one. The even, positive ground state persists into the deep near-zero spectrum;
apparent breakdowns are finite-precision or finite-cutoff artifacts, not features
of the construction.

%==============================================================================
\section{The mode budget $D_{\mathrm{nModes}}$}\label{sec:nmodes}
%==============================================================================

This is the convergence rate in $N$ --- the open problem of~\cite{CCM2025}. In
the intrinsic form it is the mode channel $\varepsilon_{\mathrm{modes}}$; as a
digit count it is the rising transition of the equilibrium-measure edge
(Section~\ref{sec:equilibrium}), climbing toward the prime ceiling.
Dividing out that ceiling (Section~\ref{sec:primes}), write the ramp ratio
$g=D_{\mathrm{nModes}}/D_{\mathrm{Primes}}\in[0,1]$. We claim
\begin{equation}\label{eq:nmodes}
  D_{\mathrm{nModes}}(N,\lambda^2)=D_{\mathrm{Primes}}(\lambda^2)\cdot
  \tanh\!\Bigl(\bigl(c\,N/(\lambda^2\ln\lambda^2)\bigr)^{q(\lambda^2)}\Bigr),
  \qquad q(\lambda^2)=\frac{2}{2\gamma+\ln\ln\lambda^2}.
\end{equation}
We build this in four steps: the stretched-exponential meaning of the ramp
(\ref{ssec:stretch}), the archimedean leading rate (\ref{ssec:arch}), the prime
correction giving $q$ (\ref{ssec:prime}), and the Shannon coordinate
(\ref{ssec:coord}).

\subsection{The ramp is a stretched exponential}\label{ssec:stretch}

\begin{theorem}[Stretched-exponential convergence and rate unification]\label{thm:stretched}
Suppose the pre-saturation mode budget satisfies $D_{\mathrm{nModes}}=(a/\ln10)\,
N^{q}(1+o(1))$ for some $a=a(\lambda^2)>0$ and exponent $q=q(\lambda^2)$, and that
the Weil-form ground eigenvalue is simple with even eigenvector. Then:
\begin{enumerate}
\item[(i)] \emph{(Stretched-exponential.)} The eigenvalue/zero error obeys
  $|\nu_1-\gamma_1|/|\gamma_1|=\exp(-aN^{q})(1+o(1))$; equivalently $q$ is the
  analyticity class of the eigenvector ($q=1$ exponential / strip-analytic,
  $0<q<1$ sub-exponential but faster than any polynomial, $q\to0$ finite Sobolev).
\item[(ii)] \emph{(Rate unification.)} Any local power-law fit
  $\mathrm{err}\sim C\,N^{-2s}$ over a window about $N$ returns the \emph{growing}
  local exponent $2s_{\mathrm{loc}}(N)=-\,d\ln\mathrm{err}/d\ln N=a\,q\,N^{q}$.
  Hence the eigenvalue-space Galerkin exponent $s(c)$ of~\cite{Groskin2026} is the
  local logarithmic slope of this same stretched-exponential --- not a fixed
  Sobolev index --- and the digit-space ramp and the eigenvalue-space rate are one
  object in two coordinates. The factor $2$ is the eigenvalue-approximation
  identity $|\varepsilon-\varepsilon_N|=O(\|\xi-\xi_N\|^2)$ for a self-adjoint
  operator with simple ground state~\cite{BabuskaOsborn1991}.
\end{enumerate}
\end{theorem}
\begin{proof}
(i) $D=-\log_{10}(\mathrm{err})\Rightarrow\mathrm{err}=10^{-D}=\exp(-aN^q)$.
(ii) $2s_{\mathrm{loc}}=-\,d\ln\mathrm{err}/d\ln N=d(aN^q)/d\ln N=aqN^q$; the
squaring is the Rayleigh-quotient second-order identity at a simple
eigenvalue~\cite{BabuskaOsborn1991}. A degenerate ground state would give only
first-order dependence and no factor $2$.
\end{proof}

Thus $q$ is not a fitting knob but the \emph{analyticity class} of the
eigenvector: $q=1$ is pure exponential decay (analytic in a strip), $0<q<1$ is
sub-exponential but faster than any polynomial, and $q\to0$ is finite Sobolev
($N^{-2s}$). The measured exponents lie just below $1$ and decrease with
$\lambda^2$ (Table~\ref{tab:q}), the signature of an eigenvector that is
strip-analytic in its archimedean part and roughened by the prime content.

The force of part~(i) is that it converts the open problem of \emph{measuring}
a convergence rate into the analytic question of \emph{identifying} the
eigenvector's regularity: the single exponent $q$ fixes the entire pre-saturation
curve. It also explains why fixed power-law (Sobolev) fits are structurally the
wrong model --- they are the $q\to0$ limit, whereas the operator sits near $q=1$
--- and hence why such fits drift and resist extrapolation
(Section~\ref{sec:discussion}).

Part~(ii) reconciles the two published pictures: the Galerkin exponent $s(c)$
measured by Groskin~\cite{Groskin2026} in eigenvalue space is the local
logarithmic slope of the present stretched-exponential, sampled at his window, so
it is not a fixed Sobolev index and must grow with $N$ --- his eigenvalue-space
exponent and the digit-space ramp are one object in two coordinates. The squaring
that links them needs the ground eigenvalue to be \emph{simple}, which is not
assumed but established as the naturally-even, isolated ground state verified
across $38$ configurations in the reproduction~\cite{AndrewsA}; were the state
degenerate, the eigenvalue error would be only first order in the eigenvector
error and the factor $2$ --- equivalently the $2s$ in Groskin's rate --- would not
appear. \emph{Status: T1 (given the pre-saturation ramp form as hypothesis).}

\subsection{The archimedean leading rate $\delta=\pi/4$}\label{ssec:arch}

The first positive zero $\nu_1$ of the Fourier--Mellin function $\widehat\xi(z)$
converges to $\gamma_1$ at a rate set by the analyticity strip of the
\emph{limiting} Fourier--Mellin function, governed by the archimedean factor of
the Weil form. That factor
is the completed real Gamma factor $\Gamma_\mathbb{R}(s)=\pi^{-s/2}\Gamma(s/2)$:
its
logarithmic derivative $\psi(\tfrac14+\tfrac{iz}{2})\sim\ln(z/2)$ is precisely
the $\ln|z|$ multiplier identified by Suzuki~\cite{Suzuki2026}, and the factor
itself decays as $|\Gamma(\tfrac14+\tfrac{iz}{2})|\sim\exp(-\pi|z|/4)$ by
Stirling. Hence the strip half-width is
\begin{equation}\label{eq:delta}
  \delta=\pi/4 .
\end{equation}

This value has a clean structural origin, one that also shows why the scales
$\tfrac12$ and $\pi/4$ both attach to the same archimedean factor. On the
critical line the archimedean log-derivative has the large-$z$ expansion
\begin{equation}\label{eq:digamma-split}
  \psi\!\left(\tfrac14+\tfrac{iz}{2}\right)=\log\frac{z}{2}
  +i\,\frac{\pi}{2}+o(1),\qquad z\to+\infty,
\end{equation}
since the argument $\tfrac14+\tfrac{iz}{2}$ runs in the imaginary
(critical-line) direction, so $\arg\bigl(\tfrac14+\tfrac{iz}{2}\bigr)\to\pi/2$.
Its real and imaginary parts carry the two scales:
\begin{itemize}
\item the \emph{real} part $\operatorname{Re}\psi(\tfrac14+\tfrac{iz}{2})\to
  \log(z/2)$ is the symbol, Suzuki's $\log|z|$ multiplier~\cite{Suzuki2026},
  whose analyticity strip has half-width $\tfrac12$ (the digamma poles at
  $z=\pm i/2$);
\item the \emph{imaginary} part $\operatorname{Im}\psi(\tfrac14+\tfrac{iz}{2})\to
  \tfrac{\pi}{2}$ sets the envelope, through
  \[
    \frac{d}{dz}\log\bigl|\Gamma_\mathbb{R}(\tfrac12+iz)\bigr|
    =-\tfrac12\operatorname{Im}\psi\!\left(\tfrac14+\tfrac{iz}{2}\right)
    \longrightarrow-\tfrac12\cdot\tfrac{\pi}{2}=-\tfrac{\pi}{4}.
  \]
\end{itemize}
The strip half-width and the envelope rate are therefore the real and imaginary
parts of the \emph{one} analytic function $\log\Gamma_\mathbb{R}(\tfrac12+iz)$,
a Hilbert (Kramers--Kronig) conjugate pair since $\Gamma_\mathbb{R}$ is zero-free
(minimum phase): its log-magnitude and phase determine each other. In one
identity,
\begin{equation}\label{eq:pi4}
  \delta=\frac{\pi}{4}=\frac12\cdot\frac{\pi}{2}=\tfrac12\arg(i),
\end{equation}
half the right angle of the critical-line direction. The two scales $\tfrac12$
and $\pi/4$ are thus not independent constants but conjugate parts of the single
factor $\Gamma_\mathbb{R}$: the strip half-width $\tfrac12$ and its harmonic
conjugate $\pi/4$.

Two points delimit what this establishes. It fixes the \emph{value}
$\delta=\pi/4$ as a property of the limiting $\Gamma_\mathbb{R}$ factor: $\delta$
is the half-width of the strip carrying the zeros of the limiting Fourier--Mellin
function, and governs the zero convergence $\nu_1\to\gamma_1$ (with the squaring
of Thm.~\ref{thm:stretched}(ii)); it is not the decay rate of the finite
eigenvector's Galerkin coefficients, a distinct band-limited, prime-modulated
quantity ($\approx0.71$ at $\lambda^2=100$), and the finite $\widehat\xi$, of
exponential type $L/2$ via its $\sin(zL/2)$ factor, is not literally the completed
$\xi$. And it is structural motivation rather than a full derivation: obtaining
$\pi/4$ from the symbol alone would require its conjugate part, which appears only
through the analytic completion $\Gamma_\mathbb{R}$, that is, through the finite
ground-state transform inheriting that factor, the realization step that remains
open. \emph{Status: T3.}

\begin{lemma}[Lattice sampling of the coefficients]\label{lem:sampling}
The Fourier--Mellin transform of the $N$-mode eigenvector has the closed
form $\widehat\xi_N(z)=2L^{-1/2}\sin(zL/2)\sum_{|j|\le N}\xi_j/(z-2\pi j/L)$
\cite{CvS2025}. On the frequency lattice $z_k=2\pi k/L$,
\[
  \widehat\xi_N(2\pi k/L)=(-1)^k\sqrt{L}\,\xi_k\ \ (|k|\le N),\qquad
  \widehat\xi_N(2\pi k/L)=0\ \ (|k|>N),
\]
so $|\xi_k|=L^{-1/2}\,\bigl|\widehat\xi_N(2\pi k/L)\bigr|$, with
$L=\ln\lambda^2$ the period of the basis $\{V_n\}$.
\end{lemma}
\begin{proof}
At $z_k=2\pi k/L$ the factor $\sin(zL/2)=(-1)^k\sin\!\bigl((z-z_k)L/2\bigr)$
has a simple zero, so every $j\ne k$ term vanishes and the $j=k$ pole is
removable:
$\lim_{z\to z_k}2L^{-1/2}(-1)^k\sin\!\bigl((z-z_k)L/2\bigr)/(z-z_k)
=(-1)^k\sqrt{L}$. For $|k|>N$ there is no $j=k$ term and the common factor
$\sin(z_kL/2)=0$.
\end{proof}

By Lemma~\ref{lem:sampling} the coefficients are the lattice samples of
$\widehat\xi_N$ at spacing $2\pi/L$ with $L=\ln\lambda^2$, so the strip
half-width~\eqref{eq:delta} gives a coefficient tail $\exp(-\delta(2\pi/L)N)$;
with the squaring of Thm.~\ref{thm:stretched}(ii) the leading digit slope is
\begin{equation}\label{eq:slope}
  \frac{\partial D_{\mathrm{nModes}}}{\partial N}
  =\frac{2\delta\,(2\pi/L)}{\ln10}=\frac{\pi^2}{\ln10\;\ln\lambda^2}
  \approx\frac{4.286}{\ln\lambda^2},
\end{equation}
a parameter-free prediction. At $\lambda^2=7$, where few primes enter and
$q\approx1$, \eqref{eq:slope} gives $2.203$ against a measured small-$N$ slope
of $2.264$ --- agreement to $3\%$ with no free parameters. At larger $\lambda^2$
the measured early slope runs above~\eqref{eq:slope} exactly as the $q<1$
concavity (below) requires.

\subsection{The prime correction: $q(\lambda^2)$}\label{ssec:prime}

The leading rate gives $q=1$. The finite Euler product
$\prod_{p\le\lambda^2}(1-p^{-s})^{-1}$ modulates the archimedean envelope by an
oscillatory prime term $S(z)=\sum_{p\le\lambda^2}p^{-1/2}\cos(z\ln p)+\cdots$
of mean-square strength
\[
  \langle S^2\rangle=\tfrac12\sum_{p\le\lambda^2}p^{-1}
  =\tfrac12(\ln\ln\lambda^2+M)+o(1)
\]
by Mertens' theorem ($M$ the Mertens constant); equivalently the full
log-product is $\sum_{p\le\lambda^2}-\ln(1-1/p)=\gamma+\ln\ln\lambda^2$. More
primes roughen the eigenvector and lower the analyticity exponent below $1$.
Modeling the reciprocal exponent as a baseline plus the prime fluctuation,
\begin{equation}\label{eq:q}
  \frac1q = c_0+c_1\ln\ln\lambda^2,\qquad
  c_0=\gamma\ \text{(archimedean baseline)},\quad
  c_1=\tfrac12\ \text{(variance prefactor)},
\end{equation}
which gives $q(\lambda^2)=2/(2\gamma+\ln\ln\lambda^2)$. The numerator $2$ is
$1/c_1$, the reciprocal of the $\tfrac12$ in the mean square; the $\ln\ln$ is
Mertens.

We test~\eqref{eq:q} against per-$\lambda^2$ fits of the full
form~\eqref{eq:nmodes} (Table~\ref{tab:q}). A linear regression
$1/q=c_0+c_1\ln\ln\lambda^2$ over $\lambda^2\in\{7,13,20,25,50,100,150\}$ yields
\[
  c_0=0.556\ (\gamma=0.577),\qquad c_1=0.478\text{--}0.508\ (\tfrac12=0.500),
  \qquad R^2=0.98 .
\]
Euler--Mascheroni $\gamma$ appears as the intercept of a regression that does
not assume it. \emph{Status: T3; the remaining lemma is a rigorous
variance-to-exponent map.}

\begin{table}[t]
\centering
\caption{Per-$\lambda^2$ fits of~\eqref{eq:nmodes} (full $\tanh$ form,
mode-limited points). The onset constant $c$ clusters near $4/\pi=1.273$; the
exponent $q$ decreases through $1$ near $\lambda^2\approx13$, matching the
archimedean crossover. $\lambda^2=200$ (sparse, $8$ points) is shown but
excluded from the regression.}
\label{tab:q}
\begin{tabular}{@{}rccccccc r@{}}
\toprule
$\lambda^2$ & 7 & 13 & 20 & 25 & 50 & 100 & 150 & 200 \\
\midrule
$q$ & 1.185 & 0.985 & 0.900 & 0.880 & 0.835 & 0.785 & 0.760 & 0.805 \\
$c$ & 1.20 & 1.25 & 1.28 & 1.29 & 1.31 & 1.28 & 1.26 & 1.31 \\
\bottomrule
\end{tabular}
\end{table}

\subsection{The Shannon coordinate}\label{ssec:coord}

\begin{lemma}[Coordinate from ceiling over slope]\label{lem:coord}
The natural coordinate of the ramp is $N/(\lambda^2\ln\lambda^2)$. The scale at
which the leading ramp meets its ceiling is
\[
  N^\ast=\frac{D_{\mathrm{Primes}}}{\partial D/\partial N}
  =\frac{(4\pi/\ln10)\lambda^2}{\pi^2/(\ln10\,\ln\lambda^2)}
  =\frac{4}{\pi}\,\lambda^2\ln\lambda^2 .
\]
\end{lemma}

The $\lambda^2$ comes from the prime ceiling (Section~\ref{sec:primes}); the
$\ln\lambda^2$ from the inverse of the archimedean slope~\eqref{eq:slope}.
The combination $\lambda^2\ln\lambda^2$ --- the time--frequency Shannon number
of the box (Landau--Pollak~\cite{LandauPollak1962};
Slepian--Pollak~\cite{SlepianPollak1961}; Landau--Widom~\cite{LandauWidom1980})
--- is therefore forced, not posited. (Measured saturation occurs at
$N/(\lambda^2\ln\lambda^2)\approx2.0$--$2.3$, larger than the linear value
$4/\pi$ because the $q<1$ ramp reaches its ceiling later than the initial-slope
line.) The form and the constant separate cleanly: writing
$N^\ast=\lambda^2\ln\lambda^2/\delta$, the grouping $\lambda^2\ln\lambda^2$ is
forced by the ceiling's $\lambda^2$ scaling and the basis interval
$L=\ln\lambda^2$ (Lemma~\ref{lem:sampling}) alone --- independently of the archimedean constant --- while
only the multiplier $4/\pi=1/\delta$ carries the rate $\delta$. That multiplier
is the data-pinned piece: the full-ramp onset constant clusters at $4/\pi$ across
the cutoffs of Table~\ref{tab:q} (mean $\approx1.27$). \emph{Status: the
coordinate form is T2 (ceiling-over-slope inputs); the multiplier $4/\pi$ is T3,
data-pinned.}

%==============================================================================
\section{The prime ceiling $D_{\mathrm{Primes}}$}\label{sec:primes}
%==============================================================================

The prime channel of the intrinsic law~\eqref{eq:main-intrinsic} is the accuracy
at infinite $N$ and $P$, set solely by which primes $p\le\lambda^2$ enter the
form. It is the relative error of the first zero at the cutoff,
\begin{equation}\label{eq:epsprimes}
  \varepsilon_{\mathrm{primes}}(\lambda^2)
  =\frac{|\nu_1(\infty,\lambda^2)-\gamma_1|}{|\gamma_1|},
\end{equation}
with $\nu_1(\infty,\lambda^2)$ the first zero at infinite basis size and fixed
$\lambda^2$: it is what remains once modes and precision are exhausted, and its
digit measurement $D_{\mathrm{Primes}}=-\log_{10}\varepsilon_{\mathrm{primes}}$ is
the ceiling the minimum~\eqref{eq:main} consumes. We fix its leading rate through
the plunge \emph{eigenvalue}, the tractable route, kept here in digit space to
reinforce the intrinsic statement with the measured quantity.

The plunge eigenvalue gives an \emph{eigenvalue floor}
$D_{\mathrm{eig}}=-\log_{10}\varepsilon^\infty(\lambda)$, the digit value of the
smallest eigenvalue. The CCM construction ties $\varepsilon^\infty$ to the
principal prolate spheroidal eigenvalue at bandwidth $\lambda$; the prolate
eigenvalue asymptotic (Connes~\cite{Connes2026}, \S6.4; Fuchs~\cite{Fuchs1964};
cf.\ Connes~\cite{ConnesProlate} and Hedenmalm~\cite{Hedenmalm2026}) gives, for
the relevant index,
\[
  1-\chi(\lambda)\sim \frac{2^{14}\sqrt{2}\,\pi^5}{3}\,
  \exp\!\bigl(-4\pi\lambda^2+\tfrac92\ln\lambda^2\bigr),
\]
so that
\begin{equation}\label{eq:eigfloor}
  D_{\mathrm{eig}}(\lambda^2)=K\lambda^2-\frac{9}{2\ln10}\ln\lambda^2-C_0,
  \qquad K=\frac{4\pi}{\ln10}=5.45750\ldots,
\end{equation}
with every constant closed-form: the leading $K=4\pi/\ln10$ is the base-$10$
rendering of the $e^{-4\pi\lambda^2}$ prolate decay rate that appears in Connes'
own framework, the $\tfrac92$ subleading-log coefficient is the Fuchs index, and
$C_0=\log_{10}\!\bigl(2^{14}\sqrt{2}\,\pi^5/3\bigr)=6.3736\ldots$ follows from the
prolate asymptotic (Groskin~\cite{Groskin2026} independently evaluates the same
\S6.4 prefactor and obtains $\log_{10}\approx6.37$, in agreement). This fixes the
leading rate of $\varepsilon_{\mathrm{primes}}$: the plunge sets $K\lambda^2$.

The eigenvalue depth is a scaffold, not the object. What the law consumes is the
zero error $\varepsilon_{\mathrm{primes}}$, which sits a few digits below the
eigenvalue floor: converting the plunge depth $\varepsilon^\infty$ into the zero
displacement $|\nu_1-\gamma_1|$, and normalizing by $|\gamma_1|$ (the
$\log_{10}|\gamma_1|\approx1.15$ digits of the relative error), gives the matching
floor
\begin{equation}\label{eq:primes}
  D_{\mathrm{Primes}}(\lambda^2)=K\lambda^2-P_m\log_{10}\lambda^2-c_m,
  \qquad P_m\approx5,\quad c_m\approx9.5.
\end{equation}
The leading $K\lambda^2$ is shared with the eigenvalue floor, the same plunge
rate, so the intrinsic content is a \emph{single} channel, not two competing
ceilings; the subleading $P_m,c_m$ are the eigenvalue-to-zero conversion together
with the $|\gamma_1|$ normalization. Equivalently the \emph{conversion gap}
$\mathrm{conv}(\lambda^2):=D_{\mathrm{eig}}-D_{\mathrm{Primes}}
=(P_m-\tfrac92)\log_{10}\lambda^2+(c_m-C_0)\approx\tfrac12\log_{10}\lambda^2+3$
is that conversion, read in digit space. The leading rate $K$ is derived (prolate
transfer); the conversion, hence $P_m$ and $c_m$, is empirical --- of it, the
$\approx1.15$-digit part of $c_m$ is the explicit $|\gamma_1|$ normalization and
the remainder is the depth-to-zero factor.

The arithmetic does not move the eigenvalue-floor slope. At the scale the
operator resolves, the prime part of the Weil form has large-$|t|$ behavior
$S(t)=\tfrac12 t^2-\gamma|t|+c_0$ (Mertens' theorems with prime-number-theorem
partial summation). On the even, mean-zero ground eigenvector each piece is inert
on the prefactor power: $\tfrac12 t^2$ acts as a rank-one \emph{odd} operator that
annihilates the even ground state; $-\gamma|t|$ equals $2\gamma K_a$, a multiple of
the metric operator $K_a=(-\Delta)^{-1}$, so it only rescales the eigenvalue and
shifts the additive constant; and $c_0$ projects out of the mean-zero subspace
(the screw-function operator split of~\cite{Suzuki2026}). None of the leading
prime terms can therefore supply a $(\lambda^2)^\Delta$ prefactor correction:
given the rate, the prefactor power is the Fuchs index $\tfrac92$ of the prolate
model (Section~\ref{sec:primes}), and the entire arithmetic contribution is
confined to the additive constant. This inertness is a statement about the
prefactor only --- the archimedean part of the form does not by itself carry the
plunge (it is $O(1)$-indefinite; see the rate-origin discussion in
Section~\ref{sec:discussion}). The empirically steeper
matching-floor slope $P_m\approx5$ is a finite-$\lambda^2$ curvature effect, not a
change in the asymptotic slope.

\emph{Status: the leading coefficient $K$ and the eigenvalue-floor slope $\tfrac92$
ride the prolate transfer (T2), and the arithmetic is provably inert on the slope;
only the matching-floor offset ($P_m$, $c_m$), the eigenvalue-to-zero conversion,
is empirical (T4).}

\paragraph{The floor is smooth, not a prime staircase.}
A natural hypothesis --- that $D_{\mathrm{Primes}}$ steps at prime entries ---
is falsified directly. With $N$ and $P$ held above floor ($N=300$, $P=200$),
sweeping $\lambda^2$ in unit steps from $4$ to $30$ gives a smooth rise of
$\approx5.3$ digits per unit $\lambda^2$; the per-unit rate at prime-power
entries is statistically indistinguishable from non-prime entries. The ceiling
is a smooth function of $\lambda^2$, consistent with the analytic
form~\eqref{eq:primes}.

\paragraph{Floor anchors.}
Floor-isolated measurements (both $N$ and $P$ above floor) give
$D_{\mathrm{Primes}}(100)=526.08$ at $(N,P)=(1000,700)$ and a saturated
$D_{\mathrm{Primes}}(200)\approx1070.4$ at $N\ge2400$. The latter sits at the
matching floor $K\cdot200-5\log_{10}200-9.5=1070.5$, confirming the two-floor
structure at the largest cutoff measured.

%==============================================================================
\section{The equilibrium-measure edge}\label{sec:equilibrium}
%==============================================================================

Sections~\ref{sec:nmodes} and~\ref{sec:primes} derived the mode ramp and the
prime ceiling by two separate classical routes: the Landau--Widom transition
width for the ramp, and the Fuchs prolate tail for the ceiling. This section
states the structural fact that ties them together --- both are features of a
single spectral edge --- so the intrinsic law~\eqref{eq:main-intrinsic} is one
edge rather than two coincident mechanisms. It adds no new constant; it is the
reason the two budgets share a form.

\subsection{The archimedean equilibrium measure}\label{ssec:eqmeasure}

Treat the limiting ordinates $\{\nu_k\}$ as a unit charge in the external field
set by the archimedean weight of the Weil form. As $\lambda^2\to\infty$ the
macroscopic distribution is the equilibrium measure $\mu_Q$ minimizing the
weighted logarithmic energy
\[
  I[\mu]=\iint\log\frac{1}{|s-t|}\,d\mu(s)\,d\mu(t)+2\int Q(t)\,d\mu(t),
  \qquad \mu\ge0,\ \textstyle\int d\mu=1,
\]
with external field $Q$ the archimedean potential whose derivative is the phase
rate $Q'(t)=\tfrac12\log(t/2\pi)$ of Section~\ref{ssec:arch}. That a limiting
spectral or zero density is the equilibrium measure of such a weighted energy
problem is classical (Saff--Totik~\cite{SaffTotik1997}; in the
orthogonal-polynomial and random-matrix setting, Deift~\cite{Deift1999}); we
invoke it here as a known tool. The Euler--Lagrange condition
$2\int\log|s-t|\,d\mu(t)=2Q(s)+\mathrm{const}$ on $\operatorname{supp}\mu$ fixes
the density; with this field it is
\begin{equation}\label{eq:rvm}
  \rho_Q(t)=\frac{1}{2\pi}\log\frac{t}{2\pi},
\end{equation}
the Riemann--von Mangoldt zero-counting density
($\int_0^T\rho_Q=\tfrac{T}{2\pi}\log\tfrac{T}{2\pi e}$). The smooth backbone of
the spectrum is thus the equilibrium measure of the archimedean field, recovered
without reference to the primes and independent of RH. \emph{Status: T3 --- the
density~\eqref{eq:rvm} is the standard archimedean main term and the variational
identification is structural; a rigorous equilibrium-measure limit for the
truncated operator is the open step.}

\subsection{One edge, two scales}\label{ssec:oneedge}

The equilibrium measure has a lower band edge: the eigenvalues descend from the
$O(1)$ bulk into the deep, positive plunge (positivity forced by the Weil form).
That edge is a single object, and the two budgets are its two scales:
\begin{itemize}
\item \emph{the transition is the mode ramp} --- the width of the crossover from
  bulk to plunge, as the resolved dimension crosses the Shannon number
  $\lambda^2\ln\lambda^2$, is the Landau--Widom log-width~\cite{LandauWidom1980},
  which is the $\tanh$ ramp of $D_{\mathrm{nModes}}$ (Section~\ref{sec:nmodes});
\item \emph{the extreme tail is the prime ceiling} --- the depth past the edge is
  the Fuchs extreme-value asymptotic~\cite{Fuchs1964}, which is
  $D_{\mathrm{Primes}}$ (Section~\ref{sec:primes}).
\end{itemize}
So the Landau--Widom width and the Fuchs tail, invoked separately above, are the
transition and the tail of the \emph{same} band edge. This is the mechanism
behind the algebra $D_{\mathrm{nModes}}=D_{\mathrm{Primes}}\cdot g$ of
Section~\ref{ssec:twoforms}: $D_{\mathrm{nModes}}$ and $D_{\mathrm{Primes}}$ are
one edge $D_{\mathrm{edge}}$ read at two resolutions, and the intrinsic
error~\eqref{eq:main-intrinsic} is correspondingly one edge --- its
finite-$N$ transition $\varepsilon_{\mathrm{modes}}$ and its extreme tail
$\varepsilon_{\mathrm{primes}}$ are the ramp and the ceiling of $\mu_Q$'s lower
edge. The unification is conceptual, not a new estimate: the microscopic
universality class of the edge (Airy, Bessel, or a Landau--Widom band edge) and
the exact deep-edge constant remain open, the latter being the same
archimedean--prime cancellation left open for $D_{\mathrm{Primes}}$ in
Section~\ref{sec:primes}. The equilibrium-measure framework itself is classical
(above); to the best of the author's knowledge, its application to the
Connes--Consani--Moscovici / Connes--van Suijlekom Weil form --- identifying the
zero-counting backbone as the archimedean equilibrium measure, and the mode ramp
and prime ceiling as the transition and extreme tail of that measure's lower
edge --- is new here. \emph{Status: T3 --- structural unification riding the
prolate transfer already used for both budgets; the framework is classical and
no new constant is claimed.}

%==============================================================================
\section{The complete formula}\label{sec:complete}
%==============================================================================

Collecting Sections~\ref{sec:wprec}--\ref{sec:equilibrium}, the law stands in the
two forms of the introduction. \emph{Intrinsically} it is the single edge of
Section~\ref{sec:equilibrium}: the relative error is the two-channel quantity
\begin{equation}\label{eq:complete-intrinsic}
  E(N,\lambda^2)=\varepsilon_{\mathrm{modes}}(N,\lambda^2)
                +\varepsilon_{\mathrm{primes}}(\lambda^2),
\end{equation}
equivalently, in digit units, the edge function
$D_{\mathrm{edge}}(N,\lambda^2)=D_{\mathrm{Primes}}(\lambda^2)\cdot
\tanh\!\bigl((c\,N/(\lambda^2\ln\lambda^2))^{q}\bigr)$, whose $N\to\infty$ ceiling
$D_{\mathrm{edge}}(\infty,\lambda^2)=D_{\mathrm{Primes}}$ is the prime floor and
whose finite-$N$ transition is the mode ramp $D_{\mathrm{nModes}}$.
\emph{Observed} at working precision $P$, it is that one edge capped by the
instrument resolution:
\begin{equation}\label{eq:complete}
  D(\lambda^2,N,P)=\min\!\Bigl(
  \underbrace{P+O(1)}_{D_{\mathrm{Wprec}}},\;
  \underbrace{D_{\mathrm{Primes}}\tanh\!\bigl((c\,N/(\lambda^2\ln\lambda^2))^{q}\bigr)}_{D_{\mathrm{nModes}}\;=\;D_{\mathrm{edge}}},\;
  \underbrace{K\lambda^2-P_m\log_{10}\lambda^2-c_m}_{D_{\mathrm{Primes}}\;=\;D_{\mathrm{edge}}(\infty)}
  \Bigr),
\end{equation}
with $K=4\pi/\ln10$ exact, $q=2/(2\gamma+\ln\ln\lambda^2)$, and $c\approx4/\pi$.
The last two entries are the one edge $D_{\mathrm{edge}}$ at finite $N$ and at
$N\to\infty$, so the operative law is $D=\min(D_{\mathrm{Wprec}},
D_{\mathrm{edge}})+O(1)$: a single spectral edge, capped by precision. This is
the Andrews CCM Convergence Law. Of its parts, the minimum
structure, the precision slope, and the stretched-exponential convergence and
rate unification (Thm.~\ref{thm:stretched}) are rigorous (T1); the Shannon coordinate and the prime
leading coefficient are derived modulo a stated lemma (T2); the archimedean rate,
the exponent $q$, and the onset constant are derived in structure and validated
to a few percent against high-precision data (T3); the prime subleading offset is
empirical (T4). Table~\ref{tab:tiers} is the precise ledger.

To the best of the author's knowledge, the complete convergence model
of~\eqref{eq:complete} is novel, meaningful, and has not been previously
published: the assembly of the three budgets into a single law, the closed-form
mode exponent $q(\lambda^2)=2/(2\gamma+\ln\ln\lambda^2)$, the Shannon coordinate
$\lambda^2\ln\lambda^2$, the stretched-exponential identification of the mode
budget, and the recognition of the mode ramp and prime ceiling as one
equilibrium-measure edge appear here for the first time. The rate it characterizes is precisely the
problem named open in~\cite{CCM2025}: Groskin~\cite{Groskin2026} measures a
Galerkin exponent but disclaims deriving it, and \'Sliwi\'nski~\cite{Sliwinski2026}
establishes only a lower bound, so no prior work supplies the functional form
given here. The novelty claimed is the law as a whole and the identification of
its mechanism; the underlying construction~\cite{CCM2025,CvS2025}, the
Paley--Wiener strip mechanism~\cite{Groskin2026}, and the classical prolate and
Mertens inputs are credited to their sources throughout.

\begin{proposition}[Convergence criterion and the prime bottleneck]\label{prop:limit}
Let $\nu_1=\nu_1(\lambda^2,N,P)$ and $D=-\log_{10}(|\nu_1-\gamma_1|/|\gamma_1|)$.
\begin{enumerate}
\item[(i)] \emph{(Prime bottleneck.)} For every fixed $\lambda^2$,
  $\sup_{N,P}D(\lambda^2,N,P)=D_{\mathrm{Primes}}(\lambda^2)<\infty$. Hence no
  sequence with bounded $\lambda^2$ has $\nu_1\to\gamma_1$: convergence to the
  zero forces $\lambda^2\to\infty$.
\item[(ii)] \emph{(Criterion.)} Since $D_{\mathrm{nModes}}\le D_{\mathrm{Primes}}$
  --- a mode-limited count cannot beat the ceiling --- $D\to\infty$ if and only if
  $D_{\mathrm{Wprec}}\to\infty$ and $D_{\mathrm{nModes}}\to\infty$; the latter
  forces $\lambda^2\to\infty$ (part~(i)), with $N$ past the Shannon scale
  $\lambda^2\ln\lambda^2$ and $P\to\infty$. The mode budget rises monotonically
  toward its ceiling as $N\to\infty$ (\cite{CCM2025}, Prop.~3.4).
\end{enumerate}
Consequently, conditional on the continuum plunge
$D_{\mathrm{Primes}}(\lambda^2)\to\infty$ (i.e.\ $\varepsilon^\infty(\lambda)\to0$),
convergence holds along any path $\lambda^2\to\infty$ with
$N\gtrsim(4/\pi)\lambda^2\ln\lambda^2$ and $P\gtrsim K\lambda^2$, at rate
$D\sim K\lambda^2$.
\end{proposition}
\begin{proof}
(i) As $N\to\infty$ the truncation converges to the continuum operator at cutoff
$\lambda$, whose first zero lies at relative distance
$10^{-D_{\mathrm{Primes}}(\lambda^2)}$ from $\gamma_1$; here
$D_{\mathrm{Primes}}(\lambda^2)<\infty$ because $\varepsilon^\infty(\lambda)>0$
(Weil positivity; \cite{Sliwinski2026}) and is $N$-independent
(\cite{Suzuki2026}, Thm 1.3), and finite working precision only lowers $D$.
Hence $D\le D_{\mathrm{Primes}}(\lambda^2)$, finite, for all $N,P$.
(ii) By Prop.~\ref{prop:min}, $D=\min(D_{\mathrm{Wprec}},D_{\mathrm{nModes}},
D_{\mathrm{Primes}})+O(1)$; since $D_{\mathrm{nModes}}\le D_{\mathrm{Primes}}$,
this is $\min(D_{\mathrm{Wprec}},D_{\mathrm{nModes}})+O(1)$, so $D\to\infty$ if and
only if $D_{\mathrm{Wprec}}\to\infty$ and $D_{\mathrm{nModes}}\to\infty$; since
$D=-\log_{10}(|\nu_1-\gamma_1|/|\gamma_1|)$, this is exactly
$\nu_1\to\gamma_1$.
\end{proof}

\noindent This refines the qualitative monotone convergence of~\cite{CCM2025}
(Prop.~3.4) into an explicit, conditional rate. The equivalence
$D\to\infty\Leftrightarrow\nu_1\to\gamma_1$ and the prime bottleneck~(i) are
rigorous (T1); the conditional rate inherits the rigor of $D_{\mathrm{Primes}}$
--- the prolate transfer (T2) and the continuum plunge it assumes.

Conceptually, the criterion isolates the prime cutoff as the sole hard
bottleneck: precision and basis size are freely improvable and impose no ceiling,
so reaching the exact zero is, at root, a statement about $\lambda^2\to\infty$ and
the growth of the prime content. The rate $D\sim K\lambda^2$ quantifies the cost
--- each additional digit demands a fixed increment $\ln10/(4\pi)$ in $\lambda^2$
--- turning the construction's qualitative convergence into a concrete, if
conditional, budget.

%==============================================================================
\section{Discussion}\label{sec:discussion}
%==============================================================================

\paragraph{Relation to Groskin~\cite{Groskin2026}.}
Groskin independently implemented the Connes--van Suijlekom matrix and measured
a Galerkin convergence exponent $s(c)$ growing with the cutoff, proposing a
Paley--Wiener analyticity-strip mechanism. The present work builds on that
mechanism: Thm.~\ref{thm:stretched}(ii) shows his $s(c)$ is the local slope of our
stretched-exponential, and the strip width is identified with the archimedean
$\Gamma_\mathbb{R}$ rate $\delta=\pi/4$. What is new here is the digit-space
formulation, the Shannon coordinate, the closed form for $q(\lambda^2)$, and the
embedding of the rate in the full three-budget minimum.

Groskin's revised analysis, extended to cutoff $\lambda^2=100$, is independent
corroboration of this two-regime structure. He finds that no single
smooth-convergence law fits his eigenvalue-space data and explicitly declines to
assert a convergence rate, reporting instead an unresolved tension between a
\emph{drifting} effective exponent and an apparent \emph{finite} convergence
limit. Both halves are accounted for here: Theorem~\ref{thm:stretched}(ii) explains
the drift --- a fixed Sobolev index cannot persist, since the local exponent
$2s_{\mathrm{loc}}(N)=aqN^{q}$ grows with $N$ --- while the prime ceiling of
Section~\ref{sec:primes} supplies the finite limit $D_{\mathrm{Primes}}(\lambda^2)$.
What looks like an irreconcilable choice within a single-rate framework is simply
the pre-saturation ramp and the ceiling of the minimum law~\eqref{eq:main}, seen
in one dataset.

\paragraph{Relation to \'Sliwi\'nski~\cite{Sliwinski2026} and Suzuki~\cite{Suzuki2026}.}
\'Sliwi\'nski proves an inverse-logarithmic lower bound on the \emph{mean}
spectral error over the whole spectrum (Thm.~3.1) --- a different regime from the
first-zero convergence here, and consistent with it, since the mean is dominated
by the unconverged high-index eigenvalues while $\nu_1\to\gamma_1$ converges
super-exponentially. Suzuki's screw-function analysis supplies the archimedean
$\ln|z|$ multiplier and the prime-frequency support that underlie
Sections~\ref{ssec:arch}--\ref{ssec:prime}.

\paragraph{Origin of the leading rate: arithmetic, not archimedean.}
The prime ceiling is \emph{modeled} by an archimedean object --- the prolate
eigenvalue of Section~\ref{sec:primes} --- yet the exponential rate is not carried
by the archimedean part of the Weil form itself. Diagonalizing the
screw-function split~\cite{Suzuki2026} at high precision, the archimedean-only
localized form is $O(1)$-indefinite, with no exponentially small eigenvalue.
Restricting it to the Sonin subspace --- the orthogonal complement of the
band-concentrated prolate modes, where Connes establishes archimedean Weil
positivity~\cite{Connes2026} --- restores positivity, but only as an $O(1)$ gap:
the smallest Sonin eigenvalue is $2.26$, $2.28$, $2.57$ at $\lambda^2=5,8,13$,
essentially flat in $\lambda^2$, against a floor $K\lambda^2=27,44,71$.
Archimedean positivity and the plunge are therefore \emph{decoupled}: the place
at infinity supplies an $O(1)$ spectral gap, while the exponentially small
$\varepsilon^\infty\sim e^{-4\pi\lambda^2}$ is produced by the cancellation
between the archimedean and prime parts of the full form. This locates the single
open analytic step behind the leading coefficient: a quantitative
archimedean--prime cancellation estimate on the even ground eigenvector --- the
precise sense in which $K$ is~T2.

\paragraph{Practical use.}
The law is a parameter-selection guide: to target $D^\ast$ digits, pick
$\lambda^2\gtrsim D^\ast\ln10/(4\pi)=D^\ast/5.46$ so the ceiling clears
$D^\ast$; pick $N$ so the $\tanh$ ramp clears $D^\ast$ (a few times the Shannon
number $\lambda^2\ln\lambda^2$); and pick $P\gtrsim D^\ast$. Increasing any one
parameter past the binding minimum yields nothing.

%==============================================================================
\section{Rigor ledger}\label{sec:rigor}
%==============================================================================

Table~\ref{tab:tiers} collects the rigor status of every component of the law,
together with the open item that would promote each to the next level.

\begin{table}[H]
\centering
\caption{Rigor status of each component. T1: rigorous. T2: derived modulo one
named lemma. T3: derived in structure and validated against data, with a map not
yet rigorous. T4: empirical.}
\label{tab:tiers}
\begin{tabular}{@{}lll@{}}
\toprule
Component & Status & Open item \\
\midrule
Minimum structure (Prop.~\ref{prop:min}) & \textbf{T1} & --- \\
$D_{\mathrm{Wprec}}=P+O(1)$, slope $1$ & \textbf{T1} & implementation offset \\
Stretched-exp convergence $+$ rate unification (Thm.~\ref{thm:stretched}) & \textbf{T1} & --- \\
Convergence criterion (Prop.~\ref{prop:limit}) & \textbf{T1} & continuum plunge $\varepsilon^\infty\!\to\!0$ (T2) \\
$D_{\mathrm{Primes}}$ leading $K=4\pi/\ln10$ & \textbf{T2} & prolate asymptotic transfer \\
$D_{\mathrm{Primes}}$ floor slope $9/2$ & \textbf{T2} & prolate transfer; primes provably inert \\
$D_{\mathrm{Primes}}$ matching offset ($P_m,c_m$) & \textbf{T4} & empirically pinned \\
Shannon coordinate $\lambda^2\ln\lambda^2$ & \textbf{T2} & ceiling/slope (Lem.~\ref{lem:coord}) \\
Equilibrium-measure edge (one edge, two scales) & \textbf{T3} & equilibrium limit; edge universality class \\
Archimedean rate $\delta=\pi/4$ & \textbf{T3} & zero-convergence from $\Gamma_\mathbb{R}$ \\
Exponent $q=2/(2\gamma+\ln\ln\lambda^2)$ & \textbf{T3} & variance$\to$exponent map \\
Onset constant $c$ ($\sim 4/\pi$) & \textbf{T3} & exact rational \\
\bottomrule
\end{tabular}
\end{table}

%==============================================================================
\section{Conclusion}\label{sec:conclusion}
%==============================================================================

We have given a single law~\eqref{eq:main} for the convergence of the CCM
zeta spectral triple, in which the minimum structure is forced, the prime
ceiling has the exact
leading coefficient $4\pi/\ln10$, and the mode budget --- the convergence rate
named open in~\cite{CCM2025} --- is the logarithm of a stretched-exponential
whose exponent is the eigenvector's analyticity class, derived end to end as
$q=2/(2\gamma+\ln\ln\lambda^2)$. The mode ramp and the prime ceiling are not two
mechanisms but one object --- the lower edge of the archimedean equilibrium
measure, whose bulk density is the Riemann--von Mangoldt law --- seen at two
scales (Section~\ref{sec:equilibrium}). Read as a convergence statement, the law says
$\nu_1\to\gamma_1$ exactly when all three budgets diverge, giving an explicit
rate $D\sim(4\pi/\ln10)\lambda^2$ along $\lambda^2\to\infty$
(Prop.~\ref{prop:limit}). Each claim is stated with an explicit rigor
level. The rigorous core (the minimum, the precision slope, and the
stretched-exponential convergence theorem~\ref{thm:stretched}, which unifies the
digit- and eigenvalue-space rates) stands on its own; the remaining lemmas --- a coefficient-decay estimate
for the $\Gamma_\mathbb{R}$-weighted form, and a variance-to-exponent map ---
would promote the archimedean rate and the exponent to theorems, and are the
natural next steps.

%==============================================================================
\section{Explicit non-claims}\label{sec:nonclaims}
%==============================================================================

The per-component rigor of the law is the subject of Table~\ref{tab:tiers} and
is not repeated here. Beyond it, we are explicit about three boundaries.

\begin{itemize}
\item \textbf{No proof of the Riemann Hypothesis.} We characterize how the
  CCM/CvS approximation converges; the link from the construction to RH ---
  continuum positivity of $QW_\lambda$ for all $\lambda$ --- is the open
  conjecture of~\cite{CCM2025,CvS2025,Connes2026}, and is neither assumed nor
  advanced here.
\item \textbf{We do not establish that the construction reproduces the zeta
  zeros.} That the finite operator's zeros converge to the true ordinates as
  $\lambda^2\to\infty$ is the construction's own open premise; we quantify the
  approach to $\gamma_1$ \emph{under} that premise (Prop.~\ref{prop:limit}) and
  do not prove the premise itself.
\item \textbf{Scope.} The law is stated and tested for the first zero $\gamma_1$
  of the Riemann zeta function. We make no claim about its constants for higher
  zeros, nor about extension to Dirichlet $L$-functions or other settings.
\end{itemize}

%==============================================================================
\section*{Statement on the use of AI tools}
%==============================================================================
In keeping with the arXiv policy on authors' use of generative AI language
tools, the author discloses that Anthropic's Claude AI model served as a
text-to-text aid in preparing the LaTeX exposition and the Rust implementation
code, to specifications set by the author. AI assistance was also used in gathering
research material; all such material was inspected and characterized by the
author. The author originated, directed, and checked every computational
experiment and mathematical claim presented here, and bears full responsibility
for the contents of this manuscript.

%==============================================================================
\section*{Acknowledgments}
%==============================================================================
The author thanks Professor Henri Moscovici (The Ohio State University) and
Professor Nigel Higson (The Pennsylvania State University) for correspondence
and encouragement, and Professor Alain Connes (Coll\`ege de France and IH\'ES)
for reviewing the author's earlier reproduction~\cite{AndrewsA} and sharing it
with fellow researcher Akiva Groskin.

%==============================================================================
\section*{Empirical testing and code availability}
%==============================================================================
The results above rest on an extensive empirical program. The author conducted
over $5{,}000$ high-precision CCM experiments to observe the operator's behavior
across the cutoff, basis-size, and working-precision axes. This testing was
critical: it provided the foundation for the derivation, ruling candidate ideas
in or out quickly and aiming the analysis at the behavior the operator actually
exhibits.

The paper source and reproduction scripts are available at
\url{https://github.com/TeamXcelerator/ccm-convergence-rate}; the shared
high-precision library (Xcelerator Toolkit) is at
\url{https://github.com/TeamXcelerator/xcelerator-toolkit}. Both are
source-available in accordance with the licensing provided with each repository.

%==============================================================================
\begin{thebibliography}{99}
%==============================================================================

\bibitem{CCM2025}
A.~Connes, C.~Consani, H.~Moscovici,
\emph{Zeta Spectral Triples}, arXiv:2511.22755 (2025).

\bibitem{CvS2025}
A.~Connes, W.~D.~van Suijlekom,
\emph{Quadratic Forms, Real Zeros and Echoes of the Spectral Action},
arXiv:2511.23257 (2025).

\bibitem{ConnesProlate}
A.~Connes,
\emph{Zeta zeros and prolate wave operators}, arXiv:2310.18423.

\bibitem{Connes2026}
A.~Connes,
\emph{The Riemann Hypothesis: Past, Present and a Letter Through Time},
arXiv:2602.04022 (2026).

\bibitem{Suzuki2026}
M.~Suzuki,
\emph{Weil's quadratic form via the screw function}, arXiv:2606.09096 (2026).

\bibitem{Hedenmalm2026}
H.~Hedenmalm,
\emph{Spectral Interpretation of Riemann Zeta Zeros}, arXiv:2606.17494 (2026).

\bibitem{Groskin2026}
A.~Groskin,
\emph{High-Precision Approximation of Riemann Zeros via the Truncated Weil
Form}, arXiv:2605.20224 (2026).

\bibitem{Sliwinski2026}
D.~\'Sliwi\'nski,
\emph{Spectral Analysis of the $D_{\log}^{(\lambda,N)}$ Operators},
arXiv:2601.12133 (2026).

\bibitem{Fuchs1964}
W.~H.~J.~Fuchs,
\emph{On the eigenvalues of an integral equation arising in the theory of
band-limited signals}, J. Math. Anal. Appl. \textbf{9} (1964), 317--330.

\bibitem{LandauWidom1980}
H.~J.~Landau, H.~Widom,
\emph{Eigenvalue distribution of time and frequency limiting},
J. Math. Anal. Appl. \textbf{77} (1980), 469--481.

\bibitem{SlepianPollak1961}
D.~Slepian, H.~O.~Pollak,
\emph{Prolate spheroidal wave functions, Fourier analysis and uncertainty --- I},
Bell Syst. Tech. J. \textbf{40} (1961), 43--63.

\bibitem{LandauPollak1962}
H.~J.~Landau, H.~O.~Pollak,
\emph{Prolate spheroidal wave functions, Fourier analysis and uncertainty ---
III: the dimension of the space of essentially time- and band-limited signals},
Bell Syst. Tech. J. \textbf{41} (1962), 1295--1336.

\bibitem{SaffTotik1997}
E.~B.~Saff, V.~Totik,
\emph{Logarithmic Potentials with External Fields},
Grundlehren der mathematischen Wissenschaften \textbf{316}, Springer, 1997.

\bibitem{Deift1999}
P.~Deift,
\emph{Orthogonal Polynomials and Random Matrices: A Riemann--Hilbert Approach},
Courant Lecture Notes in Mathematics \textbf{3}, American Mathematical Society,
1999.

\bibitem{BabuskaOsborn1991}
I.~Babu\v{s}ka, J.~Osborn,
\emph{Eigenvalue problems}, in \emph{Handbook of Numerical Analysis}, Vol.~II
(Finite Element Methods, Part~1), P.~G.~Ciarlet and J.~L.~Lions, eds.,
North-Holland, 1991, pp.~641--787.

\bibitem{AndrewsA}
R.~Andrews, Jr.,
\emph{Independent Reproduction and Convergence Analysis of the
Connes--Consani--Moscovici Zeta Spectral Triple}, Zenodo (2026),
doi:10.5281/zenodo.20427499.

\bibitem{AndrewsB}
R.~Andrews, Jr.,
\emph{Empirical Falsification of Convergence Rate Predictions for the
Connes--Consani--Moscovici Zeta Spectral Triple}, Zenodo (2026),
doi:10.5281/zenodo.20427673.

\bibitem{Arb2017}
F.~Johansson,
\emph{Arb: efficient arbitrary-precision midpoint--radius interval arithmetic},
IEEE Trans. Comput. \textbf{66} (2017), no.~8, 1281--1292.

\bibitem{Odlyzko}
A.~M.~Odlyzko,
\emph{Tables of zeros of the Riemann zeta function},
\url{http://www.dtc.umn.edu/~odlyzko/zeta_tables/}.

\end{thebibliography}

\end{document}
